\chapter{Neck Pain}

\section*{Definition}
Pain or discomfort localized to the cervical spine region, from the base of the skull to the upper shoulders, arising from musculoskeletal, neurological, or referred sources.

\section*{Pathophysiology}
Neck pain most commonly results from mechanical dysfunction: muscle strain, ligamentous sprain, or facet joint irritation. Degenerative changes (cervical spondylosis) involve disc desiccation, osteophyte formation, and facet hypertrophy, which may lead to foraminal stenosis and nerve root compression (cervical radiculopathy). Myelopathy from spinal cord compression is a serious but less common cause.

\section*{Etiology}
\begin{itemize}
  \item Mechanical neck pain (muscle strain, poor posture)
  \item Cervical spondylosis (degenerative disc disease)
  \item Cervical radiculopathy (herniated disc, foraminal stenosis)
  \item Whiplash injury (motor vehicle accident)
  \item Tension-type headache (referred to neck)
  \item Meningitis (with fever, photophobia, Kernig sign)
  \item Cervical myelopathy (spinal cord compression)
  \item Fibromyalgia or myofascial pain syndrome
  \item Rheumatoid arthritis (atlantoaxial instability)
  \item Vertebral artery dissection (rare, with neurological symptoms)
\end{itemize}

\section*{Indications for Evaluation}
Seek care if:
\begin{itemize}
  \item Severe or worsening symptoms
  \item Neck pain with fever and stiff neck (possible meningitis), neck pain after trauma, or neck pain with arm weakness/numbness
  \item Accompanied by other concerning signs
\end{itemize}

\section*{Related Symptoms}
\begin{itemize}
  \item Headache
  \item Shoulder pain
  \item Arm pain or numbness
  \item Stiffness
  \item Limited range of motion
\end{itemize}

\textbf{Note:} Always consider serious underlying conditions when evaluating persistent neck pain.

\section*{Self-Care / Home Management}
\begin{itemize}
  \item Rest and avoid activities that may aggravate the condition.
  \item Maintain adequate hydration and a balanced diet.
  \item Over-the-counter remedies may provide symptomatic relief (consult a pharmacist or doctor).
  \item Monitor symptoms and seek professional advice if they persist or worsen.
  \item Avoid self-medicating with prescription medications without medical guidance.
\end{itemize}

\section*{Diagnostic Approach}
\begin{itemize}
  \item A thorough history and physical examination are the first steps in evaluation.
  \item Laboratory tests (blood work, urinalysis) may be ordered based on clinical suspicion.
  \item Imaging studies (X-ray, ultrasound, CT, MRI) may be indicated when structural pathology is suspected.
  \item Specialist referral is arranged when the diagnosis remains uncertain or requires specific expertise.
\end{itemize}

\section*{Treatment / Management Overview}
\begin{itemize}
  \item Treatment targets the underlying cause identified through diagnostic evaluation.
  \item Supportive care and symptom management are provided while awaiting definitive therapy.
  \item Medications, lifestyle modifications, and physical therapies may be prescribed as appropriate.
  \item Follow-up ensures treatment effectiveness and allows timely adjustments to the care plan.
\end{itemize}

\clearpage